Four kinds of signal appear in this study: three synthetic ones, with which the models are probed and each of which serves a distinct purpose, and the real telemetry of the application domain. This section states what each source is and why it is used, while the appendices carry the pseudocode, the parameters and the pool composition.

\subsection*{Single-harmonic sinusoid}

It takes the form $x(t)=A\sin(2\pi f t+\phi)$, sampled at $f_s$, and is the primitive the design rests on: frequency and phase are exact, so a locked frequency produces exactly the degeneracy of \autoref{eq:token_collapse} and no background content can be blamed for a loss observed at a predicted site.

\subsection*{Light TSMixup}

A controlled probing variant of the mixup augmentation used to train Chronos\autocite{ansariChronosLearningLanguage2024}, in which the real source datasets of the original are replaced by a pool of sinusoids. It tests whether the effect survives when the tone is embedded in a multi-component signal whose composition is still fully known (\autoref{sec:appTSMixup}).

\subsection*{KernelSynth}

The Gaussian-process generator of the same training recipe\autocite{ansariChronosLearningLanguage2024}, whose output carries trends and broadband content as well as periodicity. It tests whether the effect survives in signals resembling the corpus the model was trained on rather than the probe used to measure it (\autoref{sec:appKernelSynth}).

\subsection*{Photovoltaic telemetry}

The measured series of the SolarTech Lab at Politecnico di Milano, released open access on IEEE DataPort\autocite{levaPhotovoltaicPowerWeather2020}, records six channels at one-minute resolution: module power at a $30^\circ$ tilt, ambient temperature, global horizontal and plane-of-array irradiance, wind speed and wind direction. They supply the quantities the ontology of \autoref{sec:appOntology} names and the cadence against which the dimensionless readings of \autoref{sec:tokenisation} are interpreted. The channel forecast here is the module power over 2017, of which only about a third of the minutes carry a valid reading, mostly because the plant produces and logs nothing at night. What remains is a sequence of daytime fragments: a smooth arc peaking near $200$\,W on clear days, broken on cloudy ones by short drops and peaks that come from the weather rather than from the plant, and that no forecaster can anticipate from the past alone. \autoref{sec:appSolar} shows representative days and describes how the gaps are treated.
